\cleardoublepage
\phantomsection
\pdfbookmark{Compendio}{Compendio}
\begingroup
\let\clearpage\relax
\let\cleardoublepage\relax
\chapter*{Sommario}
Il presente documento descrive il lavoro svolto durante il periodo di \textit{stage}, dal 5 maggio 2026 al 1 luglio 2026, presso l'azienda Ergon Informatica S.r.l..

\noindent
Il percorso di tirocinio è stato monitorato dal Prof. Tullio Vardanega e dal tutor aziendale Thomas Garbui e ha avuto come obiettivo principale la realizzazione di una \textit{web application} per la gestione logistica delle prenotazioni di banchine di scarico per mezzi pesanti, mediante lo sviluppo con \textit{React} e \textit{C\#}.

\noindent
Il documento è suddiviso in \textbf{4 capitoli}:
\begin{itemize}
    \item \textbf{Capitolo 1: Ergon Informatica S.r.l. - Ambito organizzativo e metodologie di lavoro}. Viene descritta la struttura organizzativa dell'azienda, le metodologie di lavoro adottate e le tecnologie utilizzate nel contesto del tirocinio.
    \item \textbf{Capitolo 2: Strategia aziendale e obiettivi del percorso di tirocinio}. Viene analizzato il ruolo del tirocinio all'interno della filosofia di innovazione, confrontando le motivazioni alla base della proposta progettuale e la rispondenza del progetto alle esigenze di mercato.
    \item \textbf{Capitolo 3: Kargo - Progetto di tirocinio}. Viene presentato il progetto di tirocinio, descrivendo il ciclo di vita del \textit{software}, le tecnologie utilizzate, le decisioni progettuali e le funzionalità implementate nell'applicativo.
    \item \textbf{Capitolo 4: Valutazione retrospettiva}. Viene effettuata un'analisi critica dei risultati ottenuti dallo \textit{stagiaire}, comparando le aspettative iniziali con i risultati finali, e misurando l'impatto del tirocinio sullo sviluppo delle competenze professionali e personali.
\end{itemize}
\newpage
\chapter*{Convenzioni Tipografiche}
Al fine di rendere più chiara la lettura del documento, sono state adottate le seguenti convenzioni tipografiche:
\begin{itemize}
    \item Il \textit{corsivo} è utilizzato per evidenziare termini in lingua straniera, titoli di opere o documenti, e per enfatizzare concetti o parole chiave.
    \item Le immagini e le figure sono numerate progressivamente e accompagnate da una didascalia descrittiva che ne chiarisce il contesto.
    \item I codici sorgente sono presentati in un font monospazio, con numerazione delle righe e indentazione coerente per facilitarne la consultazione.
    \item I termini contrassegnati dal simbolo $^{G}$ in apice sono definiti all'interno del glossario, a cui si rimanda per un approfondimento terminologico.
\end{itemize}
\endgroup
\vfill
